\section{Introduction（序論）}

\subsection{背景}

那須地域は、日本有数の高原リゾート地として長年発展してきた観光地域である。豊かな自然環境、温泉、別荘地、レジャー施設、美術館群などを背景に、多くの観光客を受け入れてきた。一方で近年では、日本全体の人口減少・高齢化・地域経済縮小の影響を受け、地域活動や個人商店の減少、後継者不足など、地方地域共通の課題も顕在化している。

特に観光地域では、「観光客向け空間」と「地域住民の日常」が分離されやすい傾向がある。観光施設は観光客を受け入れる一方、地域住民はサービス提供者として位置づけられ、双方が本質的に交わらないまま地域構造が形成される場合も少なくない。

また、大型宿泊施設は地域経済を支える重要な観光インフラである一方で、地域商店との競合対象として語られることも多い。

しかし実際には、大型宿泊施設が持つ集客力・空間・文化的資源と、地域住民が持つ趣味・文化活動・人的関係性は、本来相互に接続可能な要素でもある。

本研究では、その接続可能性に着目した。


\subsection{問題提起}

従来の観光イベントの多くは、

\begin{itemize}
\item 外部アーティスト招聘
\item 一方向型鑑賞
\item 商業型演出
\item 消費型体験
\end{itemize}

を中心として構成されている。

この形式では、観光客は「見る側」、地域住民は「提供する側」に固定されやすく、

\begin{itemize}
\item 地域参加
\item 文化継承
\item 世代交流
\end{itemize}

などが限定的となる傾向がある。

また、多くの地域イベントは「新しいものを作る」ことへ注力しやすい。しかし実際の地域には既に、

\begin{itemize}
\item 音楽
\item 人
\item 空間
\item 文化
\item 高齢者の経験
\item 観光施設
\end{itemize}

など、多くの素材が存在している。

本研究では、それら既存要素を接続し直すことで、新たな地域共存構造が形成される可能性を検証する。


\subsection{研究目的}

本研究は、ホテルサンバレー那須・観光客・地域住民という三者の関係を、「消費・労働・競合関係」ではなく、「相互創生構造」として再構築することを目的とする。

ケーススタディとして、ホテルサンバレー那須において実施された地域住民参加型クリスマス聖歌企画を対象とし、

\begin{itemize}
\item 地域住民による自主的文化活動
\item 観光客参加型歌唱
\item 歌詞カードによる参加設計
\item 主催者メッセージによる感謝
\item 聖歌背景説明による文化理解
\end{itemize}

などが、どのように観光体験へ影響したかを分析する。

また本研究では、単なるイベント成功事例としてではなく、

\begin{itemize}
\item 空間再活用
\item 地域文化継承
\item 高齢者参加
\item 観光体験設計
\item 情報設計
\end{itemize}

という複数要素が接続された社会構造として考察を行う。
